% !TeX spellcheck = en_GB
\chapter{Conclusion}
\label{ch:conclusion}



Conformal Field Theories (CFTs) are an important area of quantum theory since they can describe systems at a phase transition. In this thesis, we have studied the connection between CFTs and subfactors. Here, we have discussed several approaches to build microscopic models from Unitary Fusion Categories (UFCs) with the aim of finding evidence for a connection between the corresponding subfactor and a CFT. This study is motivated by the concrete example of the Haagerup subfactor, which is the simplest example of a subfactor for which it is not clear whether a corresponding CFT exists. The connection between subfactors and CFTs is mainly drawn by a one-dimensional lattice model, the anyon chain, which is built from the Unitary Modular Tensor Category (UMTC) that can be constructed from the subfactor. Numerically studying the behaviour of the ground state when the system size goes to infinity provides information about the corresponding CFT, given that the system exhibits a continuous quantum phase transition.

In case of the Haagerup subfactor we only know the UFCs but not the UMTC, hence we cannot directly build an anyon chain from it. However, constructing the UMTC from the UFCs is a challenging task on its own. Therefore, the simplest approach to study the connection to CFTs is to build an anyon chain directly from the UFC, even though this is not guaranteed to succeed. The most important result of this thesis, even though it is a negative one, is the insight that it is not sufficient to use the UFC in order to find a CFT via the anyon chain method: in the case of the Haagerup fusion category $\Hd$, we find no evidence for a corresponding CFT in \cref{sec:anyons}.

Therefore, it is inevitable that one needs to build the UMTC, i.e., the Drinfeld centre (or quantum double), of the fusion categories that come from the Haagerup subfactor. Note that it does not matter which of the three categories in the Morita equivalence class we use since their Drinfeld centres are the same. Here, we have discussed two techniques: First, we have introduced the annular category in \cref{ch:defects}, whose representations are connected to the simple objects of the Drinfeld centre of the category. However, this technique involves many tedious calculations, hence it has not yet been pursued for the Haagerup fusion categories. The second way of finding the Drinfeld centre, which is the more promising one, is by constructing the Levin-Wen model and calculating its excitations. Here, the biggest obstacle is to find a generalisation of the model such that it is applicable to fusion categories that do not fulfil the tetrahedral symmetry condition, which was presented in \cref{ch:LW}. However, finding the excitations for the generalised model turned out to be highly complicated due to the lack of symmetries that are present in the original model, but not in the generalised form. 

In summary, we cannot give a conclusive answer to the question we asked in the beginning of this thesis, which was whether there is a CFT that corresponds to the Haagerup subfactor. Nevertheless, through the investigations made in this thesis we have come a long way in answering this question: Firstly, they revealed which methods can \emph{not} be used for the study of this conjecture, namely building anyon chains directly from a UFC. Secondly, they have helped us develop a clear plan for the next steps that are necessary to answer this question:
	\begin{enumerate}
		\item To rule out the possibility of using UFCs instead of UMTCs completely, we can build an anyon chain from the two other fusion categories in the Morita equivalence class of $\Hd$. Since $\Hi_2$ has the same fusion rules, it will not yield any new results. However, it is possible that an anyon chain built from $\Hi_1$ (which has different fusion rules than $\Hi_2$ and $\Hd$) shows critical behaviour and hence yields the desired CFT. Even if this investigation also does not yield any evidence for a CFT, it gives new insights to building an anyon chain for a UFC with multiplicities, since the fusion rules of $\Hi_1$ are not multiplicity-free (in contrast to $\Hi_2$ and $\Hd$). This information would be valuable for the final step of the investigation, which is building an anyon chain from the Drinfeld centre, since this is a UMTC with multiplicities.
		\item To obtain a UMTC that corresponds to the Haagerup subfactor, the Drinfeld centre of the Haagerup fusion categories has to be constructed. The most promising approach here seems to be the construction of a tensor network representation of the ground state of the Levin-Wen model and computing its excitations, as explained in \cref{sec:excitations}. The first step in this approach is then to find a generalisation of the existing methods to UFCs that do not exhibit tetrahedral symmetry. This general method could then be applied to one of the Haagerup fusion categories in order to construct the Drinfeld centre.
		\item Once the Drinfeld centre for the Haagerup fusion categories is constructed, it can be used to build and numerically study an anyon chain in the same way as it is described in \cref{sec:anyons}. This would be the final step in the investigation of the connection between the Haagerup subfactor and a possible corresponding CFT.
	\end{enumerate}

Whilst the techniques presented here have been focused on the specific goal of understanding the possibility of a CFT corresponding to the Haagerup subfactor, they may also be applicable to more complex subfactors. This would help us gain a clearer picture of whether the conjectured general correspondence exists. In particular, the approach via anyon chains in \cref{sec:anyons} is applicable to all subfactors once the UMTC is known, while the UMTC itself can be obtained from the generalisation of the Levin-Wen model presented in \cref{ch:LW}.


%Apart from these concrete tasks that deal with the correspondence between subfactors and CFTs for the specific example of the Haagerup subfactor, it is an interesting question whether the techniques presented here can be used to study this correspondence for further, possibly more complex examples, or even help to find a general answer to this conjecture.